\usepackage{iftex}

%%%%%%%%%%%%%%%%%%%%%%%%%%%%%%%%%%%%%%%%%
\usepackage{amsmath,amssymb,soul}
%\usepackage[colorlinks=true, allcolors=blue]{hyperref}

\usepackage{bussproofs,caption,graphicx,float,tikz}
\usepackage{pdfpages}
\usepackage{algorithm}
\usepackage{mdframed}
\usepackage{mathpartir,enumerate}
\usepackage{algpseudocode,amsthm,tcolorbox}
\usepackage[colorinlistoftodos]{todonotes}
\usepackage{blindtext,multicol,stmaryrd}

\usepackage[english]{babel}

\usetikzlibrary{arrows.meta}
\usepackage{xy} 
\xyoption{all} 
\xyoption{knot}
\xyoption{arc}
\xyoption{import}
\xyoption{poly}

\usepackage{tikz}
\usetikzlibrary{arrows}
\usepackage{tkz-euclide}

\usepackage{lipsum}

\usepackage{array}

%%%%%%%% commands %%%%%%%%

\newcommand{\daniele}[1]{\textcolor{cyan}{#1}}
\newcommand{\mauricio}[1]{\textcolor{red}{#1}}
\usepackage{graphics} %needed for the NEW quantifier
\newcommand\new[0]{\reflectbox{\ensuremath{\mathsf{N}}}}
%The quantifier itself, a reversed sans-serif N
%(by analogy with the reversed E of `exists'
%and the reversed A of `forall')
\newcommand\New[1]{\new #1.}
%A macro, as in \New{a}\phi
\newcommand{\adis}{\mathrel{\not\approx_\alpha}}
\newcommand{\eqDis}{\overset{?}{\adis}}
\newcommand{\nfresh}{\reflectbox{\fresh}}
\newcommand{\aeq}{\mathrel{\approx_\alpha}} 
%The set of free variables.
\newcommand{\vSet}[1]{\overline{#1}}  
\newcommand{\nomPr}[1]{\mathcal{#1}}

\newcommand{\aleq}{\approx_{\alpha}}
\newcommand{\caleq}{\approx_{\alpha,\theory{C}}}
\newcommand{\ealeq}{\approx_{\alpha,\theory{E}}}

%%% nominal relations %% 
\newcommand{\enarrow}[1]{\rightsquigarrow_{\theory{#1}}}

% the New quantifier
\newcommand\NEW[0]{\reflectbox{\ensuremath{\mathsf{N}}}}
\newcommand\nw[1]{#1^{\scalebox{.4}{\NEW}}}


\newcommand{\alert}[1]{\textcolor{red}{#1}}
\newcommand{\dani}[1]{\textcolor{orange}{#1}}
\newcommand{\pink}[1]{\textcolor{magenta}{#1}}
\newcommand{\purple}[1]{\textcolor{purple}{#1}}


\newenvironment{scprooftree}[1]%
  {\gdef\scalefactor{#1}\begin{center}\proofSkipAmount \leavevmode}%
  {\scalebox{\scalefactor}{\DisplayProof}\proofSkipAmount \end{center} }

		\definecolor{bostonuniversityred}{rgb}{0.8, 0.0, 0.0}

  \newcommand{\Int}[3]{\llbracket #1 \rrbracket^{#2}_{#3}}

  \newcommand{\COREf}{{\sf CORE}_{\fix{}}}

      \newcommand{\CORE}{{\sf CORE}}

            \newcommand{\D}{\cal{D}}

            \newcommand{\F}{\mathbb{F}}


\newcommand{\pfix}{\mathcal{P}_{\fix{}}}

\newcommand{\Inst}[1]{\mathtt{Inst}(#1)}

\newcommand{\s}[3]{\mathcal{S}_{#1}(#2 \parallel #3)}

\newcommand{\fresh}{\#}
\newcommand{\atomSet}{\mathbb{A}}


\newcommand{\atm}[1]{\mathtt{Atm}(#1)}

\newcommand{\ran}[1]{\mathtt{Ran}(#1)}

\newcommand{\pair}[1]{\langle #1\rangle}

\newcommand{\U}[2]{\mathcal{U}_{#1}(#2)}

\newcommand{\CSU}[2]{\mathcal{CSU}_{#1}(#2)}

\newcommand{\muCSU}[2]{\mu\mathcal{CSU}_{#1}{#2}}

\newcommand{\V}{\mathbb{V}}

\newcommand{\id}{\mathtt{Id}}

\newcommand{\A}{\mathbb{A}}

\newcommand{\cimplies}{\Longrightarrow_\C}

\newcommand{\C}{\ensuremath\context{C}}

\newcommand{\E}{\mathtt{E}}

\newcommand{\supp}[2]{\texttt{supp}_{#1}(#2)}

\newcommand{\act}{\cdot}

\newcommand{\pr}{\mathtt{Pr}}

\newcommand{\mpr}{\mathtt{MPr}}

\newcommand{\unify}{\mathtt{unify}_{\fix{}}}

\newcommand{\cunify}{\C-\mathtt{unify}_{\fix{}}}

\newcommand{\match}{\mathtt{matching}_{\fix{}}}

\newcommand{\cmatch}{\C-\mathtt{matching}_{\fix{}}}

\newcommand{\desunify}{\mathtt{desunify}_{\fix{}}}

\newcommand{\cdesunify}{\C-\mathtt{desunify}_{\fix{}}}

\newcommand{\cons}{\mathtt{consistent}_{\fix{}}}

\newcommand{\ccons}{\C-\mathtt{consistent}_{\fix{}}}

\newcommand{\dfix}{\mathtt{D}_{\fix{}}}

%\newcommand{\inst}{\lesssim}

\newcommand{\dom}[1]{\mathtt{dom}(#1)}

\newcommand{\perm}[1]{\mathtt{perm}(#1)}

\newcommand{\lin}[1]{\overline{#1}}

\newcommand{\var}[1]{\mathtt{Var}(#1)}

\newcommand{\vran}[1]{\mathtt{VRan}(#1)}

\newcommand{\ex}[3]{#1^{\lin{#2},\lin{#3}}}

\newcommand{\impliess}{\overset{*}{\Longrightarrow}}

\newcommand{\cimpliess}{\overset{*}{\Longrightarrow}_\C}

\newcommand{\fix}[1]{\curlywedge_{#1}}

\newcommand{\falphaeq}[1]{\overset{\fix{}}{\approx}_{\alpha #1}}

\newcommand{\fCeq}[1]{\overset{\fix{}}{\approx}_{\theory{C} #1}}

\newcommand{\nf}[1]{\pair{#1}_\mathtt{nf}}
\newcommand{\nfpair}[1]{\langle #1\rangle_{nf}}
\newcommand{\pAction}[2]{#1\cdot #2}
\newcommand{\abs}[2]{[#1]#2}

\newcommand{\T}{{\tt T}}

\newcommand{\f}[1]{\ensuremath{\mathit{#1}}}

\newcommand{\theory}[1]{\ensuremath{\mathsf{#1}}}

\newcommand{\unif}{\; _?{\approx}_?\;}

\newcommand{\csecp}[2]{\textcolor{darkgray}{\sf CP}_{\theory{#1}}(\theory{#2})}

\newcommand{\barcsecp}[3]{\textcolor{darkgray}{\sf CP}_{\theory{#1}}(\theory{#2}/\theory{#3})}

\newcommand{\context}[1]{\mathbb{C}}

\newcommand{\func}[1]{{\sf #1}}

% THEOREM Environments
\theoremstyle{plain}
\newtheorem{theorem}{Theorem}[section]
\newtheorem{corollary}[theorem]{Corollary}
\newtheorem{lemma}[theorem]{Lemma}
\newtheorem{proposition}[theorem]{Proposition}
\newtheorem{result}[theorem]{Result}
\newtheorem{col}[theorem]{Corollary}
\newtheorem{conjecture}{Conjecture}[section]
\newtheorem{alemma}{Lemma}[section]
\newtheorem{atheorem}{Theorem}[section]
\newtheorem{aproposition}{Proposition}[section]



\theoremstyle{definition}
\newtheorem{definition}[theorem]{Definition}
\newtheorem{example}{Example}[section]

\theoremstyle{remark}
\newtheorem{remark}[theorem]{Remark}

